\section{Ubicación del proyecto}
\label{sec:ubicacion_proyecto}

\subsection{Localización geográfica e institucional}

La investigación se desarrolla en el Colegio Villa de las Palmas, institución educativa privada de educación media ubicada en Palmira, Valle del Cauca, Colombia. Palmira es un municipio ubicado en el norte del departamento del Valle del Cauca, aproximadamente a 30 kilómetros de Cali, con una población de cerca de 300,000 habitantes.

El Colegio Villa de las Palmas atiende estudiantes de educación básica secundaria y educación media. Al año 2026 cuenta con una matrícula de 230 estudiantes, distribuidos en los niveles de sexto a undécimo grado. La institución dispone de aulas equipadas con computadores e internet, acceso a recursos tecnológicos, y apertura de la administración para facilitar el desarrollo de investigaciones educativas en sus instalaciones.

\subsection{Población estudiantil}

La población de referencia son estudiantes de grados décimo y undécimo del colegio. Esta población fue seleccionada por varias razones:

\begin{enumerate}
\item \textbf{Madurez cognitiva:} Estudiantes en edades entre 15 y 17 años se encuentran en etapa de pensamiento formal, con capacidades de abstracción y razonamiento lógico-matemático desarrolladas.

\item \textbf{Familiaridad con trigonometría:} Es en décimo y undécimo donde se enseña trigonometría formalmente en el currículo colombiano, por lo que estos estudiantes tienen acceso a instrucción sistemática en el tema.

\item \textbf{Experiencia con tecnología:} Estudiantes de esta edad típicamente tienen experiencia previa con herramientas digitales, lo que reduce la curva de aprendizaje asociada al uso de la herramienta de IA.

\item \textbf{Disponibilidad temporal:} Es posible coordinar con la institución sesiones de trabajo regulares dentro de horario académico.
\end{enumerate}

\subsection{Cronología de la implementación}

La implementación está planeada para una duración de 8 a 10 semanas, distribuidas a lo largo de un período académico semestral. Esto permite:

\begin{itemize}
\item Sesiones de trabajo regulares con estudiantes (idealmente 2-3 sesiones por semana de 45-60 minutos cada una).
\item Tiempo suficiente para que los estudiantes completen un conjunto coherente de tareas trigonométricas.
\item Espacio para entrevistas con estudiantes seleccionados sin interrumpir demasiado el horario de clases.
\item Período de análisis de datos al final sin presión de tiempo extrema.
\end{itemize}

La cronología específica se coordina con la institución considerando calendarios académicos, períodos de evaluación, y disponibilidad de espacios físicos.

\subsection{Justificación de la selección del sitio}

Se seleccionó el Colegio Villa de las Palmas por:

\begin{enumerate}
\item \textbf{Accesibilidad:} La institución está dispuesta a permitir la investigación y facilitar acceso a estudiantes y espacios.

\item \textbf{Características pedagógicas:} Existe disposición del docente de matemáticas a participar, coordinando la implementación sin que esto interfiera significativamente con el currículo regular.

\item \textbf{Composición estudiantil:} La población es heterogénea en términos de desempeño académico en matemáticas, lo que permite documentar cómo estudiantes con distintos niveles de competencia se relacionan con la IA.

\item \textbf{Recursos tecnológicos:} El colegio tiene acceso a computadores y conectividad, lo que permite que el grupo experimental tenga acceso confiable a la herramienta de IA.

\item \textbf{Contexto socioeconómico:} Los estudiantes provienen de contextos socioeconómicos medios, lo que implica familiaridad previa con tecnología pero no ``natividad digital'' extrema que pudiera sesgar los resultados.
\end{enumerate}
